\documentclass{report}
\usepackage{amsmath,amssymb}
\setlength{\parindent}{0mm}
\setlength{\parskip}{1em}
\begin{document}
\begin{center}
\rule{6in}{1pt} \
{\large Hans De Sterck \\
{\bf Adaptive Algebraic Multigrid for Singular Value Decomposition}}

Department of Applied Mathematics \\ University of Waterloo \\ 200 University Avenue West \\ Waterloo \\ ON N2L 3G1 \\ Canada
\\
{\tt hdesterck@uwaterloo.ca}\end{center}

An adaptive (or `self-learning') algebraic multigrid (AMG) method is
described for computing a few of the largest or smallest singular values
of a rectangular matrix, and their associated singular vectors. The
proposed algorithm combines and extends two existing AMG approaches for
symmetric positive definite eigenvalue problems. The method consists of
two multilevel phases. In the first, multiplicative phase (setup phase),
tentative singular triplets are calculated along with a multigrid
hierarchy of interpolation operators that approximately fit the tentative
singular vectors in a collective and self-learning manner, using
multiplicative update formulas and bootstrap AMG interpolation. In the
second, additive phase (solve phase), the tentative singular triplets are
improved up to the desired accuracy by using an additive correction
scheme with fixed interpolation operators, combined with a Ritz update. A
suitable generalization of the singular value decomposition is formulated
that applies to the coarse levels of the multilevel cycles. The combined
multiplicative-additive approach results in an AMG method for extremal
singular triplets that combines two desirable properties: it allows for
high-accuracy convergence when desired, and it is flexible enough to deal
efficiently with a variety of problems due to its self-learning
properties. Numerical tests show that the algorithm converges to high
accuracy in a modest number of iterations for singular value and
eigenvalue problems from different application areas. Extension of the
approach to canonical tensor decomposition is also briefly discussed.


\end{document}
